%%%%%%%%%%%%%%%%%%%%%%%%%%%%%%%%%%%%%%%%%%%%%%%%%% 
% Picture with arrows - mathrm
%%%%%%%%%%%%%%%%%%%%%%%%%%%%%%%%%%%%%%%%%%%%%%%%%% 

\usetikzlibrary{arrows}

% Define how TiKZ will draw the nodes
\tikzset{mathterm/.style={draw=white,fill=white,rectangle,anchor=base}}
\tikzstyle{every picture}+=[remember picture]
\everymath{\displaystyle}

\makeatletter

%%%%%%%%%%%%%%%%%%%%%%%%%%%%%%%%%%%%%%%%%%%%%%%%%% 
% Designate a term in a math environment to point to
% Syntax: \mathterm[node label]{some math}
%%%%%%%%%%%%%%%%%%%%%%%%%%%%%%%%%%%%%%%%%%%%%%%%%% 

\newcommand\mathterm[2][]{%
   \tikz [baseline] { \node [mathterm] (#1) {$#2$}; }}

%%%%%%%%%%%%%%%%%%%%%%%%%%%%%%%%%%%%%%%%%%%%%%%%%% 
% A command to draw an arrow from the current 
% position to a labelled math term
% Default color=black, default arrow head=stealth
% Syntax: \indicate[color]{term to point to}[path options]
%%%%%%%%%%%%%%%%%%%%%%%%%%%%%%%%%%%%%%%%%%%%%%%%%% 

\newcommand\indicate[2][black]{%
   \tikz [baseline] \node [inner sep=0pt,anchor=base] (i#2) {\vphantom|};
   \@ifnextchar[{\@indicateopts{#1}{#2}}{\@indicatenoopts{#1}{#2}}}
\def\@indicatenoopts#1#2{%
   {\color{#1} \tikz[overlay] \path[line width=1pt,draw=#1,-stealth] (i#2) edge (#2);}}
\def\@indicateopts#1#2[#3]{%
   {\color{#1} \tikz[overlay] \path[line width=1pt,draw=#1,-stealth] (i#2) [#3] edge (#2);}}

\makeatother

%%%%%%%%%%%%%%%%%%%%%%%%%%%%%%%%%%%%%%%%%%%%%%%%%% 
% Footnote for title
%%%%%%%%%%%%%%%%%%%%%%%%%%%%%%%%%%%%%%%%%%%%%%%%%% 

\let\ACMmaketitle=\maketitle
\renewcommand{\maketitle}{\begingroup\let\footnote=\thanks \ACMmaketitle\endgroup}

\def\tang{\ThisStyle{\abovebaseline[0pt]{\scalebox{-1}{$\SavedStyle\perp$}}}}

%%%%%%%%%%%%%%%%%%%%%%%%%%%%%%%%%%%%%%%%%%%%%%%%%% 
%
% Insert coding
%
% Example: 
% \lstinputlisting[language=Octave, numbers=none]{code/fsine.m}
%%%%%%%%%%%%%%%%%%%%%%%%%%%%%%%%%%%%%%%%%%%%%%%%%% 

\definecolor{codegreen}{rgb}{0,0.6,0}
\definecolor{codegray}{rgb}{0.5,0.5,0.5}
\definecolor{codepurple}{rgb}{0.58,0,0.82}
\definecolor{backcolor}{rgb}{0.95,0.95,0.92}
\definecolor{mygreen}{RGB}{28,172,0}

\lstdefinestyle{mystyle}{
    %backgroundcolor=\color{backcolor},   
    commentstyle=\color{mygreen},
    keywordstyle=\color{blue},
    numberstyle= \tiny\color{codegray},
    stringstyle=\color{codepurple},
    basicstyle=\ttfamily\footnotesize,
    %frame=L,
    %frame = lines,
    %frame = single,
    breakatwhitespace=false,         
    breaklines=true,                 
    captionpos=b,                    
    keepspaces=true,
    %numbers=left, 
    %numbersep=7pt,                  
    showspaces=false,                
    showstringspaces=false,
    showtabs=false,                  
    tabsize=2,
    % xleftmargin=5pt,
    % framexleftmargin=10pt,
    % framextopmargin=4pt,
    % framexbottommargin=4pt,
    %emph=[1]{for,end,break},emphstyle=[1]\color{blue},
    %emph=[2]{length,gca}, emphstyle=[2]\color{black}
}
 
\lstset{style=mystyle}

%%%%%%%%%%%%%%%%%%%%%%%%%%%%%%%%%%%%%%%%%%%%%%%%%%%%%
\newcounter{mycode}
\colorlet{colexam}{green!75!black}
\tcbset{
  base/.style={
    breakable,
    enhanced,
    frame engine = path,
    sharp corners,
    title,
    attach boxed title to top left,
    boxed title style={
            size=minimal, top=0pt, left=0pt
            },
    coltitle=black,
    fonttitle=\bfseries,
  }
}
%
\newtcolorbox[use counter=mycode]{code}{%
  base,
  boxed title style={
  overlay={
      \node[anchor=north east,inner ysep=0pt,align=right,text width=2in] 
        at ([yshift=-0.55ex]frame.north west) {\hfill {In [\thetcbcounter]:}};
  }},
  borderline west={3pt}{0pt}{colexam},
  borderline north={1pt}{0pt}{gray!30},
  borderline south={1pt}{0pt}{gray!30},
  borderline east={1pt}{0pt}{gray!30},
  frame hidden,
  colback = gray!5,
}

%%%%%%%%%%%%%%%%%%%%%%%%%%%%%%%%%%%%%%%%%%%%%%%%%% 
% Volume in Fluid
%%%%%%%%%%%%%%%%%%%%%%%%%%%%%%%%%%%%%%%%%%%%%%%%%% 

\makeatletter
\DeclareRobustCommand{\volume}{\text{\volumedash}V}
\newcommand{\volumedash}{%
  \makebox[0pt][l]{%
    \ooalign{\hfil\hphantom{$\m@th V$}\hfil\cr\kern0.08em\textbf{--}\hfil\cr}%
  }%
}

%%%%%%%%%%%%%%%%%%%%%%%%%%%%%%%%%%%%%%%%%%%%%%%%%% 
% CREATING NICE RED BOX
% https://tex.stackexchange.com/questions/719491/how-can-i-get-the-title-in-tcolorbox-partly-framed-and-not-framed
%%%%%%%%%%%%%%%%%%%%%%%%%%%%%%%%%%%%%%%%%%%%%%%%%% 

% Defining counter:
\newcounter{defcounter}
\setcounter{defcounter}{0}
\renewcommand{\thedefcounter}{\arabic{defcounter}}

\newtcolorbox{definition}[1][]{%
    enhanced,
    before skip=5mm, after skip=5mm, 
    colback=white, colframe=red, % Colors
    rounded corners, arc=3mm,
    drop fuzzy shadow, % Shadow
    fonttitle=\bfseries,
    title={#1}
}


%%%%%%%%%%%%%%%%%%%%%%%%%%%%%%%%%%%%%%%%%%%%%%%%%% 
% quick derivative shorthand: \dxy{u}{x} → ∂_x u
%%%%%%%%%%%%%%%%%%%%%%%%%%%%%%%%%%%%%%%%%%%%%%%%%% 
\newcommand{\dxy}[2]{\partial_{#2}\,#1}
